\chapter{Realisation et implementation}
\label{chap:realisation}

\section{Environnement technique}

La plateforme est developpee en Java 17 sur Spring Boot, avec PostgreSQL comme
systeme de persistance et Ollama comme moteur d'inference local pour l'assistant
conversationnel. L'ensemble est conteneurise avec Docker et orchestre par Docker
Compose (\cref{fig:deployment}), ce qui permet de reproduire l'environnement complet
-- base de donnees, backend, moteur LLM et telechargement automatique du modele --
sur n'importe quelle machine disposant de Docker.

\section{Pipeline d'ingestion des donnees}
\label{sec:pipeline-ingestion}

\subsection{Reception et securisation du fichier}

Un fichier depose par l'utilisateur passe par trois controles avant d'etre accepte
(\cref{lst:store}) : detection du type MIME reel par Apache Tika (independamment de
l'extension declaree, pour eviter qu'un fichier malveillant deguise en
\texttt{.xml} ne soit traite), calcul d'une empreinte cryptographique, et rejet si
cette empreinte correspond a un fichier deja traite. Ce dernier point repond
directement au probleme identifie en phase de cadrage (\emph{app\_notes\_dxc.txt},
section 4) : sans deduplication, un utilisateur qui re-uploade par erreur le meme
export declencherait un nouveau traitement par lots inutile.

\begin{lstlisting}[caption={Deduplication par empreinte avant stockage (TempFileStorageServiceimpl)},label={lst:store}]
String digest = computeHash(file);
if (fileRepository.existsByHash(digest)) {
    throw new FileProcessedException();
}
fileRepository.deleteByName(type.name().toLowerCase());
fileRepository.save(new File(UUID.randomUUID(), null, file.getSize(),
        detectMimeType(file), Instant.now(), type.name().toLowerCase(),
        digest, false));
\end{lstlisting}

\subsection{Traitement par lots (Spring Batch)}

Le traitement effectif est confie a Spring Batch. Chaque type de source (Qualys,
CMDB) possede son propre \emph{job}, compose d'une seule etape orientee bloc
(\emph{chunk}) de $100$ elements : un lecteur \texttt{StaxEventItemReader} parcourt le
XML en flux (sans le charger integralement en memoire) en isolant les fragments
correspondant a l'element racine attendu (\texttt{HOST} pour Qualys, \texttt{row}
pour la CMDB), un \emph{processor} transforme chaque DTO en entite via un
\emph{mapper} MapStruct, et un \texttt{JpaItemWriter} persiste le lot.

\begin{lstlisting}[caption={Etape Spring Batch d'ingestion CMDB (BatchConfig)},label={lst:batch}]
@Bean
public Step cmdbStep(JobRepository jobRepository,
        MultiResourceItemReader<CmdbDto> reader,
        PlatformTransactionManager tx,
        JpaItemWriter<CmdbCi> writer,
        CmdbProcessor processor) {
    return new StepBuilder(jobRepository)
            .<CmdbDto, CmdbCi>chunk(100)
            .transactionManager(tx)
            .reader(reader)
            .processor(processor)
            .writer(writer)
            .build();
}
\end{lstlisting}

Chaque nouvel import est traite comme la verite complete et actuelle de sa source :
avant de lancer le job, \texttt{JobService.wipeForReingest} supprime l'integralite
des lignes precedentes du type concerne (\texttt{Host} pour Qualys, en cascade sur
ses \texttt{VulnInfo} ; \texttt{CmdbCi} pour la CMDB). Ce choix (documente en
\cref{sec:perimetre}) evite qu'un actif retire du parc ne reste indefiniment visible
dans le tableau de bord, au prix de l'absence d'historique detaille des CI entre deux
imports -- une limite assumee et reprise en \cref{sec:perspectives}.

Le lancement d'un job est protege par un verrou applicatif
(\texttt{ReentrantLock}) : les \emph{readers} StAX sont resolus tardivement via les
parametres du job (liaison par contexte de thread), et deux jobs executes en
parallele pourraient faire lire a l'un d'eux le motif de fichiers de l'autre. Le
verrou garantit qu'un seul job d'ingestion s'execute a la fois.

\subsection{Recalcul du risque : parallelisme et coherence transactionnelle}

Une fois l'ingestion terminee, le recalcul de risque
(\texttt{RiskCalculationService.\allowbreak recomputeAll}) traite chaque hote reconcilie de
maniere independante et parallele : pour chaque actif, le calcul de l'ACS et celui
du score QDS pondere sont soumis a un pool de threads dedie au calcul
(\texttt{CALC\_EXECUTOR}), puis combines des que les deux resultats sont disponibles.

\begin{lstlisting}[caption={Parallelisation du calcul par actif (RiskCalculationService)},label={lst:parallel}]
CompletableFuture<Double> acsFuture =
    CompletableFuture.supplyAsync(() -> computeAcs(cmdbCi, roleWeights, envWeights), calcExecutor);
CompletableFuture<Double> weightedQdsFuture =
    CompletableFuture.supplyAsync(() -> computeWeightedQds(host, config), calcExecutor);

return acsFuture.thenCombine(weightedQdsFuture,
    (acs, weightedQds) -> applyRiskScore(host, cmdbCi, config, acs, weightedQds));
\end{lstlisting}

Le declenchement du recalcul n'est jamais appele directement par le controleur
d'upload : il reagit a un evenement applicatif
(\texttt{IngestionCompletedEvent}), consomme par un \emph{listener}
transactionnel configure en phase \texttt{AFTER\_COMMIT}. Ce choix, deja illustre au
\cref{fig:seq-ingestion}, evite une classe de bug subtile mais reelle : sans cette
precaution, le listener pourrait demarrer sa propre transaction de lecture avant que
la transaction d'ecriture des nouvelles donnees ne soit effectivement validee, et
recalculer le risque sur un etat de base partiellement ecrit.

\section{Tableau de bord temps reel}

Le tableau de bord combine deux mecanismes : un point d'entree REST classique
(\texttt{GET} \url{/api/v1/dashboard/metrics}) qui calcule et renvoie l'ensemble des
agregats a l'ouverture de la page, et un flux Server-Sent Events
(\texttt{GET} \url{/api/v1/dashboard/stream}) auquel le navigateur reste connecte. A
chaque fin de recalcul de risque, le meme calcul d'agregats est effectue une seule
fois cote serveur et pousse a tous les clients connectes, qui redessinent leurs
graphiques Chart.js sans recharger la page (\cref{fig:seq-ingestion}, etapes
13 a 16). Ce choix evite le double ecueil du \emph{polling} (requetes repetees
inutiles) et du rechargement manuel.

\section{Assistant conversationnel (RAG)}

Ce module reste, au moment de la redaction de ce rapport, au stade du developpement
et de la conception : il s'agit d'un prototype fonctionnel demontrable, qui n'a pas
encore fait l'objet d'un deploiement en environnement de recette (\emph{staging}) ni
d'une campagne de validation comparable a celle menee sur le moteur de risque
(\cref{chap:tests}).

L'assistant s'appuie sur Langchain4j (\texttt{@AiService}) et sur un modele
executes localement via Ollama (\texttt{mistral:7b}), plutot que sur une API LLM
externe payante -- un choix coherent avec la sensibilite des donnees manipulees
(inventaire de vulnerabilites et de criticites metier). L'architecture retenue
privilegie le \emph{tool calling} : plutot que d'indexer les donnees dans une base
vectorielle et de faire de la recherche par similarite, l'agent dispose d'outils
qui executent des requetes reelles sur la base au moment ou la question l'exige,
garantissant que les reponses portent sur l'etat courant des donnees et non sur un
instantane potentiellement perime.

\begin{lstlisting}[caption={Definition de l'agent conversationnel (AiAgentService)},label={lst:aiagent}]
@AiService
public interface AiAgentService {
    @SystemMessage("""
            You are an assistant with access to tools that query a live database.
            When a question can be answered by a tool, you MUST invoke the tool.
            Never describe, explain or write out the call - actually perform it.
            Never output code.""")
    Flux<String> chat(String question);
}
\end{lstlisting}

\section{Conteneurisation et deploiement}

Le fichier \texttt{docker-compose.yml} definit quatre services : la base
PostgreSQL (avec \emph{healthcheck} pour garantir que le backend ne demarre qu'une
fois la base disponible), le backend Spring Boot, le moteur Ollama (avec reservation
GPU pour l'inference), et un conteneur ephemere \texttt{ollama-init} charge de
telecharger automatiquement le modele au premier demarrage. Un volume nomme
persiste les fichiers uploades (\texttt{/app/uploads}) independamment du cycle de
vie du conteneur applicatif.
